% -----------------------------------------------------------------------------
% Chapter 16: Synthetic & Benchmark Data Strategy
% -----------------------------------------------------------------------------
\chapter{Synthetic \& Benchmark Data Strategy}
\label{chap:data_strategy}

\section{Dual-Track Data Methodology}
A foundational challenge in deploying machine learning within underserved credit environments is the \textit{cold-start data paradox}: commercial banks cannot originate loans without historical performance data, yet historical default data cannot be gathered without originating loans. Rather than masking this constraint, CreditTech adopted a rigorous, dual-track data engineering strategy:

\begin{enumerate}
  \item \textbf{Track 1: Real-World Benchmark Fusion ($N = 22,000$):}\\
  To train and rigorously evaluate model discrimination, calibration, and hyperparameter stability on real borrower repayment behaviors, we ingested and harmonized two gold-standard international credit datasets: the \textit{Home Credit Default Risk} benchmark~\cite{home_credit_2018} (12,000 records) and the \textit{Give Me Some Credit (GMSC)} benchmark~\cite{give_me_credit_2011} (10,000 records). Together, this benchmark cohort comprises 22,000 real borrower observations with 2,098 empirical defaults (9.54\% prevalence).
  
  \item \textbf{Track 2: Controlled Synthetic Early-Adopter Cohort:}\\
  To stress-test persona boundaries, role-based authorization, geographic village partitioning, and edge-case handling (officer overrides, active grievances, missing rails) within the running web application, a controlled, deterministic synthetic dataset was engineered via \texttt{scripts/seed\_early\_adopters.py}.
\end{enumerate}

Figure~\ref{fig:data_strategy_schematic} illustrates this dual-track data architecture.

\begin{figure}[H]
\centering
\begin{tikzpicture}[
  scale=0.92, every node/.style={transform shape},
  node distance=1.0cm and 0.6cm, auto
]

  % Core Strategy Root
  \node[baseblock, fill=primary!12, draw=primary, text width=5.8cm, minimum height=1.2cm] (root) {
    \textbf{CreditTech Data Architecture}\\[2pt]
    \scriptsize Dual-Track Benchmark \& Synthetic Pilot Strategy
  };

  % Track 1: Real Benchmarks
  \node[serviceblock, below left=1.0cm and 0.3cm of root, fill=accent!12, draw=accent, text width=5.6cm, minimum height=1.7cm] (real_track) {
    \textbf{Track 1: Real Benchmarks}\\[1pt]
    \scriptsize $N=22,000$ Real Records (Home Credit + GMSC)\\
    \scriptsize 2,098 Defaults (9.54\%) $\bullet$ 5-Fold GroupKFold CV
  };

  % Track 2: Synthetic Pilot
  \node[serviceblock, below right=1.0cm and 0.3cm of root, fill=warn!12, draw=warn, text width=5.6cm, minimum height=1.7cm] (synth_track) {
    \textbf{Track 2: Synthetic Pilot Cohort}\\[1pt]
    \scriptsize Deterministic Seeding Pipeline ($N=100$)\\
    \scriptsize 15 Village Clusters $\bullet$ Overrides \& Grievances
  };

  % Output Artifacts (Clean Rectangular Baseblocks)
  \node[baseblock, below=1.0cm of real_track, text width=5.2cm, minimum height=1.4cm, fill=accent!10, draw=accent] (mrm_artifacts) {
    \textbf{Model Registry \& Figures}\\[2pt]
    \scriptsize 5 Tuned Model Benchmarks\\
    \scriptsize 13 Publication Decision Plots
  };

  \node[baseblock, below=1.0cm of synth_track, text width=5.2cm, minimum height=1.4cm, fill=warn!10, draw=warn] (db_artifacts) {
    \textbf{Relational DB \& Web Portal}\\[2pt]
    \scriptsize Operational SQLite / PostgreSQL\\
    \scriptsize 169 Passing Verification Tests
  };

  % Flows
  \draw[flowline] (root.south) -| (real_track.north);
  \draw[flowline] (root.south) -| (synth_track.north);
  \draw[flowline] (real_track) -- (mrm_artifacts);
  \draw[flowline] (synth_track) -- (db_artifacts);

\end{tikzpicture}
\caption{Research Block Diagram: Dual-Track Real Benchmark and Synthetic Pilot Data Architecture.}
\label{fig:data_strategy_schematic}
\end{figure}

\section{Real Benchmark Fusion \& Rural Feature Mapping}
The 22,000 real borrower records were mapped to the 21-feature 5-Cs rural credit framework using empirical covariance relationships:
\begin{itemize}
  \item \textbf{Character (SHG Peer Thrift):} In Home Credit, short-term debt repayment punctuality was mapped to SHG internal loan repayment rates ($[0.70, 1.00]$) and meeting attendance ratios.
  \item \textbf{Capacity (Agricultural Cashflows):} Revolving debt ratios and monthly income metrics from GMSC were calibrated to rural PM-Kisan government transfers and seasonal crop revenues.
  \item \textbf{Conditions (Sentinel-2 Remote Sensing):} Geographic cluster assignments mapped historical default variance to seasonal NDVI vegetative vigor drops and monsoon rainfall indices.
\end{itemize}

\section{Controlled Synthetic Early-Adopter Seeding}
The synthetic seed pipeline (\texttt{scripts/seed\_early\_adopters.py}) deterministically instantiates specific borrowers across two pilot clusters:
\begin{itemize}
  \item \textbf{Cluster Alpha (Chandauli District --- Canal Irrigated):}
    \begin{itemize}
      \item \textit{Radhika Devi} (\texttt{BORROWER-A1}): Pristine SHG thrift, high crop vigor (NDVI: 0.74), Score: 745 (\texttt{APPROVE}).
      \item \textit{Sita Kumari} (\texttt{BORROWER-A2}): Marginal landholder, seasonal payment lag, Score: 630 (\texttt{MANUAL REVIEW}).
    \end{itemize}
  \item \textbf{Cluster Beta (Mirzapur District --- Rain-Fed):}
    \begin{itemize}
      \item \textit{Ramu Patel} (\texttt{BORROWER-B1}): Low NDVI due to local drought, Score: 510 (\texttt{REJECT}), but approved via an officer dissent override due to a new village solar borewell.
      \item \textit{Meena Verma} (\texttt{BORROWER-B2}): Disputed electricity utility arrears, Score: 580, active open grievance under 30-day statutory SLA.
    \end{itemize}
\end{itemize}
This controlled cohort guarantees that every complex user journey, dissent override, grievance loop, and authorization boundary can be reliably and repeatedly demonstrated.
